\documentclass[10pt,a4paper]{amsart}
\usepackage[margin=19mm]{geometry}
\usepackage{amsmath,amssymb,mathtools}
\usepackage[scr=rsfs,cal=euler]{mathalfa}
\usepackage{xfrac}
\usepackage[unicode,hidelinks,bookmarks=false]{hyperref}
\newcommand{\ie}{i.e.\ }
\newcommand{\cf}{cf.\ }
\newcommand{\resp}{respectively\ }
\newcommand{\Subsection}{\subsection}
\providecommand{\nobreakdash}{\nobreak\hbox{-}}
\setlength{\emergencystretch}{2em}
\multlinegap0pt
\usepackage{bm}
\usepackage{bbm}
\usepackage{dsfont}
\usepackage{xspace}
\usepackage{cancel}
\usepackage{mathtools}
\usepackage{amsthm}
\makeatletter
\@ifundefined{thm}{\newtheorem{thm}{Thm}}{}
\@ifundefined{cor}{\newtheorem{cor}[thm]{Corollary}}{}
\@ifundefined{correspondence}{\newtheorem{correspondence}[thm]{Correspondence}}{}
\@ifundefined{lem}{\newtheorem{lem}[thm]{Lemma}}{}
\@ifundefined{prop}{\newtheorem{prop}[thm]{Proposition}}{}
\makeatother
\DeclareMathOperator{\Aut}{Aut}
\DeclareMathOperator{\Fix}{Fix}
\DeclareMathOperator{\id}{id}
\newcommand{\cE}{\mathcal{E}}
\newcommand{\s}{\sigma}
\newcommand{\tC}{\widetilde{C}}
\title[Pryms of $\mathbb{Z}\_3\times\mathbb{Z}\_3$ coverings of genus]{Pryms of $\mathbb{Z}\_3\times\mathbb{Z}\_3$ coverings of genus 2 curves}
\author{Pawe{\l} Bor\'owka, Anatoli Shatsila}
\date{Source: arXiv:2503.23041v2 (CC BY 4.0)}
\begin{document}
\maketitle

\noindent\emph{Source and licence.} Pryms of $\mathbb{Z}_3\times\mathbb{Z}_3$ coverings of genus 2 curves. Version arXiv:2503.23041v2 (CC BY 4.0), distributed under the Creative Commons Attribution 4.0 International licence, as recorded in the arXiv licence metadata for this exact version and shown on the abstract page https://arxiv.org/abs/2503.23041v2. Full rights notice: licenses/arxiv-2503.23041-CC-BY-4.0.txt.\par\smallskip

\noindent\textit{Editorial scope.} Counted unit: the Prop whose source label is \texttt{global} in the article (source0.tex lines 530--533, source label \texttt{global}), delivered together with the article's own complete original proof of it (source0.tex lines 535--554, one \texttt{proof} environment of 20 lines). No mathematics is written, repaired or re-derived: statement and proof are the article's own text line for line. Required context retained uncounted: decprym (lines 317--321); kernel (lines 331--334); degree (lines 388--397); correspondence (lines 508--516). Every definition, display and label of those lines is retained unchanged. Omitted from this package and not redistributed: 1--316, 322--330, 335--387, 398--507, 517--529, 534--534, 555--720. Nothing inside a retained excerpt is omitted or altered: every retained body is delivered exactly as the article's lines are written, so no omission receipt of any kind accompanies this package. Citations: the retained excerpts carry no citation command at all, so no bibliography entry of the article is transcribed and no reference is redistributed. Reference adaptation: none was needed and none was made; every reference inside the delivered unit resolves inside it, so no unresolved reference or citation is produced. Printed slips kept literal: no printed word, symbol, hypothesis or number of the counted statement or of its proof was altered. Layout and class: the article's own class is \texttt{amsart} with packages including babel, geometry, amsmath, amssymb, graphicx, hyperref, mathtools, amsfonts, amsthm, tikz-cd, polski, comment, caption; the wrapper is the lane's pinned \texttt{amsart} editorial wrapper at 10pt on A4, which restates the theorem-like environments the retained text needs and adds the article's own macro definitions that the retained text needs (\texttt{\textbackslash Aut}, \texttt{\textbackslash Fix}, \texttt{\textbackslash cE}, \texttt{\textbackslash id}, \texttt{\textbackslash s}, \texttt{\textbackslash tC}), with extra wrapper packages (bm, bbm, dsfont, xspace, cancel, mathtools). The delivered numbering is the unit's own. Assets: no retained span contains a figure environment, a graphics-inclusion command, a PDF-page inclusion, a table or a TikZ picture; no image file is copied into this package, the package declares no asset dependency and no image file is redistributed with it. Licence: version arXiv:2503.23041v2 of this article is distributed under the Creative Commons Attribution 4.0 International licence, as recorded in the arXiv licence metadata for this exact version and shown on the abstract page https://arxiv.org/abs/2503.23041v2. A full rights notice is delivered at licenses/arxiv-2503.23041-CC-BY-4.0.txt. Characters: every retained excerpt is pure ASCII, so the delivered engine, XeLaTeX with its default Latin Modern text font, reports no missing character.
\begin{cor}
\label{decprym}
    The isotypical decomposition of $P=P(f)$ is as follows:
    $$P = k_{\sigma}^*P(h_{\sigma}) \boxplus k_{\tau}^*P(h_{\tau}) \boxplus k_{\sigma\tau}^*P(h_{\sigma\tau}) \boxplus  k_{\sigma^2\tau}^*P(h_{\sigma^2\tau}).$$
\end{cor}

\begin{lem}\label{trivialkernel}
\label{kernel}
    The (restricted) kernel $\ker k_{\sigma}^*|_{P(h_{\sigma})}$ is trivial in the non-isotropic case, and has cardinality $3$ in the isotropic case. In particular, in the non-isotropic case, we have $P_\s\simeq P(h_{\sigma})$.
\end{lem}

\begin{lem}
\label{degree}
     Let $$\Phi: f^*JH \times P_{\s} \times P_{\tau} \times P_{\s\tau} \times P_{\s^2\tau} \to J\tC$$ and $$\mu: P_{\s} \times P_{\tau} \times P_{\s\tau} \times P_{\s^2\tau} \to P(f)$$ be the addition maps. Then $$\deg \Phi = \begin{cases}
         3^{14}, \text{ if $f$ is non-isotropic}, \\
         3^{10}, \text{ if $f$ is isotropic.}
     \end{cases}$$ and $$\deg \mu = \begin{cases}
         3^{10}, \text{ if $f$ is non-isotropic}, \\
         3^{6}, \text{ if $f$ is isotropic.}
     \end{cases}$$
\end{lem}

\begin{correspondence}
\label{correspondence}
\begin{align*}
    M &\ \ \ \ \ \ \cE\\
    \text{an involution} &\longrightarrow \text{identity on 1 elliptic curve in each small Prym}\\
    \text{3 involutions that are identity on it} &\longleftarrow \text{an elliptic curve}\\
    \text{a unique involution that is identity on them} &\longleftarrow \text{any 2 elliptic curves (not in the same triple)}
    \end{align*}
\end{correspondence}

\begin{prop}
\label{global}
    Let $q$ be a group-compatible involution on $P(f)$. Then $q \in M$.
\end{prop}

\begin{proof} 
    
    Let $q \in \Aut(P(f))$ be a group-compatible involution. 
    Since $q$ is $\langle \sigma \rangle$-compatible and $\langle \tau \rangle$-compatible without loss of generality we can assume  that $q$ acts as the identity on $E_{\sigma, j} \subset P_{\sigma}$ and $E_{\tau, j} \subset P_{\tau}$. In particular $q=j$ on $P_{\sigma}+P_{\tau}$ (where $+$ denotes an algebraic sum of abelian subvarieties).
    
           
   By the correspondence \ref{correspondence}, for any two elliptic curves in  $P_{\sigma\tau}$ and $P_{\sigma^2\tau}$ there is a unique involution $j\sigma^a\tau^b$ which acts as the identity on both curves. Hence $q=j\s^a\tau^b$ on $P_{\sigma\tau}+P_{\s^2\tau}$ for some $a,b$.
   
      Thus, from the isotypical decomposition of $P(f)$ (Corollary \ref{decprym}) we have the following equality for $x_1 \in P_{\sigma},\ x_2 \in P_{\tau},\ x_3 \in P_{\sigma\tau},\ x_4 \in P_{\sigma^2\tau}$: $$q(x_1 + x_2 + x_3 + x_4) = j(x_1 + x_2) + j\sigma^a\tau^b(x_3 + x_4).$$ 
      
      Assume that $q \notin M$, hence $\sigma^a\tau^b \neq \id$. We will show that this assumption contradicts Lemma \ref{degree} by bounding the cardinality of the kernel of the addition map $\mu$ from above. Take a tuple $(x_1, x_2, x_3, x_4) \in P_{\sigma} \times P_{\tau} \times P_{\sigma\tau}\times P_{\sigma^2\tau}$ such that $x_1 + x_2 + x_3 + x_4 = 0$. Then $0 = j(x_1 + x_2) + j\sigma^a\tau^b(-x_1 - x_2),$ so that $x_1 + x_2 = \sigma^a\tau^b(x_1 + x_2)$. Analogously, $x_3 + x_4 = \sigma^a\tau^b(x_3 + x_4)$.
    Since one of the small Pryms is fixed by $\sigma^a\tau^b$, without loss of generality   we can assume $\sigma^a\tau^b = \sigma$, which implies that $\sigma(x_2) = x_2$.
    Therefore, for any $(x_1, x_2, x_3, x_4)$ in the kernel of the addition map $$\mu: P_{\sigma} \times P_{\tau} \times P_{\sigma\tau}\times P_{\sigma^2\tau} \to P(f)$$ we have $x_2 \in \Fix(\sigma, \tau)$.  Note that $x_3 = \sigma(\sigma^2(x_3))$ and, since $\sigma^2(x_3) \in P_{\tau}$ we have $\sigma^2(x_3) \in \Fix(\s)$ which implies that $x_3 \in \Fix(\sigma)$. Similarly, we get $\{x_1, x_2, x_3, x_4\} \in \Fix(\sigma, \tau)$. Now, $$P_{\sigma} \cap \Fix(\sigma, \tau) = P_{\sigma} \cap f^*JH = k_{\sigma}^*P(h_{\sigma}) \cap k_{\s}^*h_{\s}^*JH = k_{\s}^*(P(h_\s) \cap h_{\s}^*JH),$$ where the last equality follows from the fact that $\ker k_{\sigma}^{*}|_{P(h_{\s})} \subset P(h_\s) \cap h_{\s}^*JH$ in both isotropic and non-isotropic cases. Note that $P(h_\s) \cap h_{\s}^*JH \simeq \mathbb{Z}_3 \times \mathbb{Z}_3$ since the restricted polarisation from $JC_{\sigma}$ to $P(h_{\sigma})$ is of type $(1,3)$. It follows from Lemma \ref{kernel} that in the non-isotropic case $P_{\sigma} \cap \Fix(\sigma, \tau) \simeq \mathbb{Z}_3 \times \mathbb{Z}_3$ and in the isotropic case $P_{\sigma} \cap \Fix(\sigma, \tau) \simeq \mathbb{Z}_3$. 
    
    It follows from the above that in the non-isotropic (resp. isotropic) case the cardinality of $\ker \mu$ is at most $3^8$ (resp. $3^4$), which contradicts Lemma \ref{degree}. Therefore, $q$ must be in $M$.
   

   
    
    \end{proof}

\end{document}
